% =======================================
% MEETING 9
% ---------------------------------------
% DATE   - September 18, 2026
% FORMAT - F2F
% =======================================
% TOPIC  - Tensor product
% =======================================
\subsubsection{Tensor Product}
	\label{subsec:tensor product}
	The point of a tensor is this: A tensor is a function (or a \textit{machine}, if you will),
	that accepts covectors and vectors with $k$ and $l$ slots. After all the slots are filled,
	then the tensor spits out a real number. However, as long as those slots aren't filled yet,
	it becomes a tensor of lower rank, or lower type.

	Ordering matters between the sets of vectors and covectors. You
	can permute between them, but permuting within the group is not allowed. 

	Recall the identity we set out earlier. We have the following 
	coordinate basis for covectors and vectors:
	\begin{align*}
		\overline{e}_\mu = \p_\mu && \underline{e}^\mu = dx^\mu
	\end{align*}

	We can find the components of an arbitrary vector or covector likeso:
	\[
		v^\mu = dx^\mu[v] \equiv v[dx^\mu]
	\]
	and
	\[
		\omega_\mu = \omega[\p_\mu]\equiv \p_\mu[\omega]
	\]
	This is a sort of product of operation which is symmetric, hence it is called the 
	interior product.
	Coacting, or \textit{multiplying} rank 1 tensors together
	results on another tensor (recall: $a\in\R$ is just a rank 0 tensor).

	We need to define the notion of a tensor product, which
	allows us to combine tensors of different rank to produce a 
	new tensor. 

	\begin{definition}{Tensor Product}
	    If $\varphi$ and $\vartheta$ are covectors, or $(0,1)$ tensors, then
		their tensor product $\varphi\otimes\vartheta$ is a $(0,2)$ tensor, such that, for $x,y\in(1,0)$:
		\begin{equation}
			\label{eq:tensor-product-defn}
			(\varphi\otimes\vartheta)[x,y] = \varphi[x]\vartheta[y]
		\end{equation}

		More generally, if $T$ has type $(k, l)$ and $S$ has type $(m, n)$,
		then $T \otimes S$ has type $(k + m, l + n)$, defined by
		\begin{align*}
			&(T \otimes S)(\omega_1, \ldots, \omega_k, \omega'_1, \ldots, \omega'_m,
			v_1, \ldots, v_l, v'_1, \ldots, v'_n)
			=\\ &T(\omega_1, \ldots, \omega_k, v_1, \ldots, v_l)
			  \,S(\omega'_1, \ldots, \omega'_m, v'_1, \ldots, v'_n).
		\end{align*}
		This is also known as an outer product.
	\end{definition}

	\paragraph{Example}

	Say we have one-forms $dx^1$ and $dx^2$. If we let them act on a pair 
	of vetors $x$ and $y$, we have
	\begin{align*}
		(dx^1\otimes dx^2) [x,y] &= \dd{x}^1[x]\dd{x^2}[y] \\
		&=x^1 y^2
	\end{align*}
	Those are individual components.
	
	It should be clear from the definition that the outer product is
	not commutative.
	\[
		\varphi\otimes\vartheta 
		\neq 
		\vartheta\otimes\varphi
	\]
	For example, if we let it act on the vectors $x,y$, then
	\begin{align*}
		\varphi\otimes\vartheta[x,y] &=
		\varphi[x]\vartheta[y]\\
		\vartheta\otimes\varphi[x,y]&=
		\vartheta[x]\varphi[y]\\
		\varphi[x]\vartheta[y]&\neq
		\vartheta[x]\varphi[y]
	\end{align*}
	\paragraph{Properties} The tensor product must satisfy the following 
	properties
	\begin{enumerate}[(1)]
		\item \text{Bilinearity}

			For scalars $a, b$ and tensors
			  $\varphi, \theta, \omega$ of compatible type,
			  \[
				  \varphi \otimes (a\theta + b\omega)
				  = a\,(\varphi \otimes \theta) + b\,(\varphi \otimes \omega),
			  \]
			  \[
				  (a\varphi + b\theta) \otimes \omega
				  = a\,(\varphi \otimes \omega) + b\,(\theta \otimes \omega).
			  \]
		\item \text{Associativity}

			\[(\varphi \otimes \theta) \otimes \omega
			= \varphi \otimes (\theta \otimes \omega)\]
		\item \text{Multiplicative scaling}

			\[(a\varphi) \otimes \theta
			   = a\,(\varphi \otimes \theta)
		   = \varphi \otimes (a\theta)\]
		\item \text{Noncommutativity}
			
			In general, for tensors $\phi$ and $\theta$ acting on arbitary tensors,
			\[\varphi \otimes \theta \neq \theta \otimes \varphi\]
	\end{enumerate}
	\begin{proof}
	   Exercise to the student. 
	\end{proof}

\subsubsection{Basis Tensors}
	Like vectors and covectors, tensors which are defined as a linear combination of 
	scalars by basis vectors or basis covectors, we must find a way
	to generalize what a basis is for higher order tensors.

	Consider a tensor $g \in \binom{0}{2}$, acting on two vectors
	$x, y \in \binom{1}{0}$, such that
	\[
		g[x, y] \mapsto \alpha \in \R.
	\]
	We can express $x$ and $y$ in terms of a basis:
	\[
		x = x^\mu \p_\mu,
		\qquad
		y = y^\nu \p_\nu.
	\]
	Thus, we can express $g$ as
	\[
		g[x, y]
		= g\bqty{x^\mu \p_\mu,\; y^\nu \p_\nu}.
	\]
	Since $x^\mu$ and $y^\nu$ are just numbers, i.e., constants, we
	can pull them out:
	\[
		g[x, y]
		= x^\mu y^\nu\, g[\p_\mu, \p_\nu].
	\]
	Notice that $g$ is accepting two vectors. Hence, the entire
	expression is just a product of real numbers. Let us define
	\[
		g[\p_\mu, \p_\nu] = g_{\mu\nu}.
	\]
	Thus we can rearrange it to
	\[
		g[x, y]
		= \tensor{g}{_{\mu\nu}}\, x^\mu y^\nu.
	\]
	Recall that $x^\mu = dx^\mu[x]$ and $y^\nu = dx^\nu[y]$. Then
	\[
		g[x, y]
		= \tensor{g}{_{\mu\nu}}\,
		  dx^\mu[x]\, dx^\nu[y].
	\]
	From the definition of the tensor product,
	\[
		dx^\mu[x]\, dx^\nu[y]
		= (dx^\mu \otimes dx^\nu)[x, y].
	\]
	We can then represent the entire expression as
	\[
		g[x, y]
		= \tensor{g}{_{\mu\nu}}\,
		  (dx^\mu \otimes dx^\nu)[x, y].
	\]
	However, the vectors $x$ and $y$ are arbitrary. We can remove
	them, hence
	\[
		g = g_{\mu\nu}\, dx^\mu \otimes dx^\nu.
	\]
	Notice here that we have represented a tensor
	$g \in \binom{0}{2}$ as a linear combination of
	$\binom{0}{2}$ tensors $dx^\mu \otimes dx^\nu$, with
	components $g_{\mu\nu}$.

	Recall what we have done in the previous sections. A vector, for
	instance, can be represented by a product of components and a
	basis, and similarly for covectors:
	\[
		v = v^\mu \p_\mu,
		\qquad
		u = u_\nu\, dx^\nu.
	\]
	We have arrived at a similar expression. Therefore, we can
	conclude that $dx^\mu \otimes dx^\nu$ must be a basis. This is
	what we call a \textit{basis tensor}.

	In the same logic, if you have a tensor $m \in \binom{2}{0}$
	acting on two covectors in $\binom{0}{1}$, then we arrive at the
	same conclusion:
	\[
		m = m^{\mu\nu}\, \p_\mu \otimes \p_\nu,
	\]
	where $m^{\mu\nu} = m[dx^\mu, dx^\nu]$. Therefore
	$\p_\mu \otimes \p_\nu$ is also a basis tensor.

	\begin{definition}{Basis Tensor Representation}
		In general, we can represent an arbitrary $(k, l)$ tensor
		$T$ as
		\[
			T
			=
			\tensor{T}{^{a_1 a_2 \cdots a_k}_{b_1 b_2 \cdots b_l}}
			\p_{a_1} \otimes \p_{a_2} \otimes \cdots \otimes \p_{a_k}
			\otimes
			dx^{b_1} \otimes dx^{b_2} \otimes \cdots \otimes dx^{b_l}.
		\]
		where
		\[
			\p_{a_1} \otimes \cdots \otimes \p_{a_k}
			\otimes
			dx^{b_1} \otimes \cdots \otimes dx^{b_l}.
		\]
		is a $(k, l)$ basis tensor, and
		\[
			\tensor{T}{^{a_1 \cdots a_k}_{b_1 \cdots b_l}}
		\]
		are the $(k, l)$ tensor components.
	\end{definition}

	As with vectors and covectors, basis tensors must follow the
	property
	\[
		e^\mu[e_\nu] = e_\nu[e^\mu] = \tensor{\delta}{^\mu_\nu}.
	\]
	An equivalent statement for a general basis tensor is that the
	$(k, l)$ basis tensor acts on the corresponding basis vectors and
	covectors to give products of Kronecker deltas:
	\begin{align*}
		&\pqty{
			\p_{a_1} \otimes \cdots \otimes \p_{a_l}
		      \otimes 
			dx^{b_1} \otimes \cdots \otimes dx^{b_k}
		}
		\bqty{
			dx^{c_1}, \ldots, dx^{c_l}
			,\;
			p_{d_1}, \ldots, \p_{d_k}
		}\\
		&=
		\tensor{\delta}{^{c_1}_{a_1}} \cdots \tensor{\delta}{^{c_k}_{a_k}}\,
		\tensor{\delta}{^{b_1}_{d_1}} \cdots \tensor{\delta}{^{b_l}_{d_l}}.
	\end{align*}
	This is the natural generalization of $e^\mu[e_\nu] =
	\delta^\mu{}_\nu$ to higher-rank basis tensors: the covector
	slots and vector slots each pick off one Kronecker delta, and
	the full result is the product. It is this property that makes
	the representation
	\[
		T
		=
		\tensor{T}{^{a_1 \cdots a_k}_{b_1 \cdots b_l}}
		\,
		\p_{a_1} \otimes \cdots \otimes \p_{a_k}
		\otimes
		dx^{b_1} \otimes \cdots \otimes dx^{b_l}
	\]
	unique, and it is what allows the components to be extracted by
	the rule
	\[
		\tensor{T}{^{a_1 \cdots a_k}_{b_1 \cdots b_l}}
		=
		T\bqty{dx^{a_1}, \ldots, dx^{a_k},\;
		       \p_{b_1}, \ldots, \p_{b_l}}.
	\]
	In other words, the $(a_1, \ldots, a_k;\, b_1, \ldots, b_l)$ component
	of $T$ is obtained by feeding $T$ the corresponding basis
	covectors and basis vectors.
	
	This sort of notation is relatively uncommon in mathematics, but has been 
	universally accepted in physics.

	\paragraph{Example} Consider a $(2,1)$ tensor represented in tensor basis form
	\[
		\kappa = \tensor{\kappa}{^{\mu\nu}_{\sigma}} \p_\mu \otimes \p_\nu \otimes dx^\sigma
	\]
	If we want to find the $\tensor{T}{^{\alpha\beta}_\gamma}$-th component, we feed it the 
	basis tensors $dx^\alpha, dx^\beta,\text{ and }\p_\gamma$. 
	\begin{align*}
		\kappa
		[dx^\alpha,dx^\beta, \p_\gamma] 
		&= 
		\tensor{\kappa}{^{\mu\nu}_{\sigma}} 
		\p_\mu \otimes \p_\nu \otimes dx^\sigma
		[dx^\alpha,dx^\beta, \p_\gamma] \\
		&=
		\tensor{\kappa}{^{\mu\nu}_{\sigma}} 
		\p_\mu[dx^\alpha]
		\p_\nu[dx^\beta]
		dx^\sigma[\p_\gamma]\\
		&=
		\tensor{\kappa}{^{\mu\nu}_{\sigma}} 
		\tensor{\delta}{^\alpha_\mu}
		\tensor{\delta}{^\beta_\nu}
		\tensor{\delta}{^\sigma_\gamma}\\
		&=
		\tensor{\kappa}{^{\alpha\nu}_{\sigma}} 
		\tensor{\delta}{^\beta_\nu}
		\tensor{\delta}{^\sigma_\gamma}\\
		&=
		\tensor{\kappa}{^{\alpha\beta}_{\sigma}} 
		\tensor{\delta}{^\sigma_\gamma}
		\\
		&=
		\tensor{\kappa}{^{\alpha\beta}_{\gamma}} 
	\end{align*}
	
	Thus, we can see that the same contraction operation applies here as it 
	does for vectors and covectors.
\subsubsection{Symmetric Tensor Product}
	A tensor product, in general, is noncommutative.
	\[
		\p_\mu \otimes \p_\nu \neq \p_\nu \otimes \p_\mu
	\]
	However, there are important tensors for which interchanging the
	two tensor factors leaves the tensor unchanged. Such tensors are
	called \textit{symmetric tensors}. These tensors have a notion of commutativity.	

	Consider a tensor $g\in(2,0)$. 
	\[
		g = 
		g_{\mu\nu} dx^\mu \otimes dx^\nu
	\]
	In a 2D manifold (to keep the math simple), its explicit decomposition
	would be
	\begin{align*}
		g = 
		g_{11} \dd{x}^1 \otimes dx^1 +
		g_{12} \dd{x}^1 \otimes dx^2 +
		g_{21} \dd{x}^2 \otimes dx^1 +
		g_{22} \dd{x}^2 \otimes dx^2
	\end{align*}
	Now, suppose the tensor is symmetric, i.e.,
	\[
		g_{\alpha\beta} = g_{\beta\alpha} 
	\]
	This is the same sense of a matrix being symmetric. There
	are many cases where the tensor you are working with are symmetric,
	for example, in GR, the general metric tensor is symmetric. The 
	stress tensor $\sigma_{ij}$ and strain tensor $\varepsilon_{ij}$ is symmetric, the inertia tensor $I$, 
	and so on.

	In that case, we can write the explicit decomposition as
	\[
		g = 
		g_{11} \dd{x}^1 \otimes dx^1 +
		g_{12} (\dd{x}^1 \otimes dx^2 +
		\dd{x}^2 \otimes dx^1) +
		g_{22} \dd{x}^2 \otimes dx^2
	\]
	Notice, however, that all the tensor elements are now symmetric ($g_{11}=g_{11}, g_{22}=g_{22}$, and $g_{12}=g_{21}$, by construction).

	We have a sense of a symmetric product here, which we can use. It suggests that,
	for symmetric tensors, it is convenient to suppress the explicit tensor product
	and introduce an implicit symmetric product. Let us define the following

	\begin{definition}{Symmetric Tensor Product}
		For a symmetric tensor, i.e.,
		\[
			dx^{\mu}\otimes dx^\nu = dx^\nu \otimes dx^\mu
		\]
		We define a symmetric tensor product $dx^\mu dx^\nu$ such that
		\begin{align*}
			dx^{\mu}dx^{\nu} &= 
			\frac{1}{2}\pqty{
				dx^{\mu}\otimes dx^\nu + dx^\nu \otimes dx^\mu
			}\\
							 &= dx^\nu dx^\mu
		\end{align*} 
		This is also known as the symmetrized tensor product.
	\end{definition}

	Now that we can do that, we can rewrite $g$ as 
	\[
		g = g_{11}(dx^1)^2 + 2g_{12} dx^1 dx^2 + g_{22} (dx^2)^2
	\]
	Notice, if you interpret this physically, that 
	this is just an infinitesimal square of arclength.

	Now that we can write it in an easier geometric way, we just collapse the expression back to
	the general tensor product, with the understanding that the product of the differentials
	represents an appropriate symmetric tensor product.
	\[
		g = g_{\mu\nu} dx^\mu dx^\nu
	\]
	Notice that this is just a normal covector operation.
\subsubsection{Tensor Transformation Law}
	We can examine
	how the tensor components $g_{\mu\nu}$ transform by performing a change 
	of basis, from $x^\mu$ to $y^\sigma$. 

	Let $x^\mu = x^\mu(y^\sigma)$. Recall the change of basis formulation defined in the last chapter. That is,
	\[
		dx^\mu = \pdv{x^\mu}{y^\sigma} dy^\sigma
	\]
	We can then write the expression as
	\begin{align*}
		g 
		&= g_{\mu\nu} dx^\mu dx^\nu\\
		&= g_{\mu\nu} 
		\pqty{
			\pdv{x^\mu}{y^\sigma} 
			dy^\sigma 
		}
		\pqty{
			\pdv{x^\nu}{y^\rho} 
			dy^\rho 
		}
	\end{align*}
	The Jacobians are just constants, we can pull them out
	\[
		g = 
			\pdv{x^\mu}{y^\sigma} 
			\pdv{x^\nu}{y^\rho} 
			\,
		g_{\mu\nu} 
			dy^\sigma
			dy^\rho 
	\]
	On the other hand, since $g$ is the same tensor regardless of the 
	coordinate system, we may express it directly in terms of the new coordinates
	as
	\[
		g = g'_{\sigma\rho} dy^\sigma dy^\rho
	\]
	Comparing the two expressions, we now have a tensor transformation rule. We can equate 
	$g'_{\sigma\rho}$ with the coefficients on the previosu form, and we find the behavior of
	how the components of a $(0,2)$ tensor transforms under a change of basis.

	\begin{theorem}{Symmetric Tensor Transformation}
		Under a change of basis from $x$ to $y$, a symmetric $(0,2)$ tensor's components
		transform such that
		\begin{equation}
			\label{eq:symmetric-tensor-transformation}
			g'_{\sigma\rho} 
			=
			\pdv{x^\mu}{y^\sigma} 
			\pdv{x^\nu}{y^\rho} 
			\,
			g_{\mu\nu} 
		\end{equation}    
	\end{theorem}
	
	In particular,
	the metric components transform covariantly because each lower index
	contributes one factor of the Jacobian.

	This all hinges on the assumption that the tensor product is symmetric, and the tensor components are symmetric.
	The tensor product is hence, implicit.
	\paragraph{Tensor Field}
	If we assign a tensor of type $(k,l)$ at every point on a manifold $\mcal M$, we then obtain  a
	\textbf{tensor field} of type $(k,l)$. A tensor field $T$ is said to be \textbf{smooth} if 
	\[
		T(\omega^1,\ldots,\omega^k, v_1,\ldots, v_l)\in \mcal{F}_{\mcal M}
	\]
	for all smooth dual vector fields $\omega^1,\ldots,\omega^k$, and smooth vector
	fields $v_1,\ldots, v_l$.

	From the operation above, we can generalize how the tensor components 
	transform for higher order tensors. The transformation relation for the components of a tensor of type $(k,l)$ in a two coordinate systems in general is as follows
	\begin{theorem}{Tensor Transformation Law}
		For a tensor $T\in(k,l)$, with initial basis $x$, being transformed to another 
		basis $x'$, the general relation for the components is as follows:
		\begin{equation}
			\label{eq:tensor-transformation-law}
			{T'}\indices{^{{\mu_1}\ldots{\mu_k}}_{{\nu_1}\ldots{\nu_l}}}
			=
			\pdv{x'^{\mu_1}}{x^{\rho_1}} 
			\cdots
			\pdv{x'^{\mu_k}}{x^{\rho_k}} 
			\pdv{x^{\sigma_1}}{x'^{\nu_k}} 
			\cdots
			\pdv{x^{\sigma_l}}{x'^{\nu_l}} 
			\tensor{T}{^{{\rho_1}\ldots{\rho_k}}_{{\sigma_1}\ldots{\sigma_l}}}
		\end{equation}
	    
	\end{theorem}
	

\subsubsection{Contraction Operation}
	We need to define what a contraction is. It is an operation that we have encountered before, but 
	have no concrete definition yet. As such,
	
	\begin{definition}{Contraction}
		A contraction operator $C$ is a map from tensors of type $(k,l)$ to tensor
		of type $(k-1,l-1)$.
		\[
			C:\binom{k}{l}\to \binom{k-1}{l-1}
		\]
		It reduces the rank of a tensor by 1 upstairs and 1 downstairs. 
		
		For a tensor $T\in(k,l)$, a contraction $\tensor{C}{^i_j}$ acts
		such that
		\[
			\tensor{C}{^i_j}T:= \tensor{T}{^\mu_\mu} = \sum_\mu T(\cdots, e^\mu,e_\mu)
		\]
	\end{definition}
	If a tensor $T$ has a contravariant index and a covariant index
	which are contracted, then the contraction operation $C$ acts on
	the tensor by summing over the repeated index. It is equivalent to feeding 
	a tensor the $i$-th basis vector and $j$-th dual basis.
	If 
	\[
		T = \tensor{T}{^{\mu_1,\cdots,\mu_k}_{\nu_1,\cdots,\nu_l}}
		\p_{\mu_1}\otimes\cdots\otimes\p_{\mu_k}\otimes\dd{x^{\nu_1}}\otimes\cdots\dd{x^{\nu_l}}
	\]
	Then a contraction $\tensor{C}{^i_j}$ acts like
	\[
		\tensor{C}{^i_j}T = T[\p_{\mu_i},\dd{x^{\nu_l}}]
	\]
	Thus, CT is a new tensor with both its contravariant and covariant
	rank reduced by 1
	
	This is an important property, since it allows us to get the sense of a \textit{trace}
	of a tensor, which allows us to simplify tensor operations. 

	This is quite an abstract notion, so let me show you some examples.
	\begin{enumerate}[(1)]
		\item 
			Suppose $T$ has components
			\[
				T
				=
				T^{a_1\cdots \mu\cdots a_k}{}_{b_1\cdots \mu\cdots b_l}
				\,
				\p_{a_1}\otimes\cdots\otimes \p_\mu\otimes \cdots \otimes\p_{a_k}
				\otimes
				dx^{b_1}\otimes\cdots\otimes dx^\mu\otimes\cdots\otimes dx^{b_l},
			\]
			where the repeated index $\mu$ appears once upstairs and once
			downstairs. The contraction sums over all possible values of
			$\mu$:
			\[
				CT
				=
				\sum_{\mu}
				T^{a_1\cdots \mu\cdots a_k}{}_{b_1\cdots \mu\cdots b_l}
				\,
				\p_{a_1}\otimes\cdots
				\otimes
				dx^{b_1}\otimes\cdots.
			\]
			The contracted index $\mu$ then disappears, leaving a tensor of
			type $(k-1,l-1)$.
		\item
			Consider a tensor $T\in(1,1)$ with basis
			$dx^\mu$ and $\p_\mu$:
			\[
				T = T^\mu{}_\nu\,\p_\mu\otimes dx^\nu.
			\]
			When we perform a contraction, the vector and covector basis
			are evaluated against one another:
			\[
				dx^\nu[\p_\mu]=\delta^\nu_\mu.
			\]
			Thus,
			\[
				CT
				=
				T^\mu{}_\nu\,\delta^\nu_\mu
				=
				T^\mu{}_\mu,
			\]
			which is a scalar, of type $(0,0)$. This is the trace of the $(1,1)$ tensor.
		\item
			Consider a tensor 
			\[
				T = \tensor{T}{^{\mu_1,\cdots,\mu_k}_{\nu_1,\cdots,\nu_l}}
				\p_{\mu_1}\otimes\cdots\otimes\p_{\mu_k}\otimes\dd{x^{\nu_1}}\otimes\cdots\dd{x^{\nu_l}}
			\]
			A contraction $\tensor{C}{^4_5}$, where we specify that we are contracting the 4th covector
			and fifth vector, signifies that
			\[
				\tensor{C}{^4_5} = T[dx^{\mu_4}, \p_{\mu_5}]
			\]
			Which gives us a new tensor. 
		\item
			In the current notation, tensor contraction is quite cumbersome. 
			In abstract index notation, which we will discuss in the next chapter (so just read ahead of that 
			to understand what we discuss here), we 
			can show a contraction rather easily.

			For a tensor 
			\[
				\tensor{T}{^{a_1,\ldots,a_k}_{b_1,\ldots,b_l}}
			\]
			A contraction is simply replacing the 4th covector slot and 5th covector slot with the same index.
			\[
				C^4_5 T = \tensor{T}{^{a_1,\ldots,c ,\ldots,a_k}_{b_1,\ldots,c,\ldots,b_l}}
			\]

			For example
			\[
				\tensor{F}{^{ab}_c} u^c \equiv F\otimes u
			\]
			You can verify that the components are $\tensor{T}{^{\alpha\beta}_\mu} u^\mu$. We then sum over the $\mu$, and
			find a new tensor of lower rank
			\[
				\tensor{F}{^{ab}}
			\]
	\end{enumerate}

\subsubsection{Abstract Index Notation}
	When it comes to tensor notation, there are several competing schools of thought
	regarding how it would be represented. We can show these 
	different notation with the following
	\begin{enumerate}[(1)]
		\item Index--free notation

			To illustrate this, let's show the following. Consider a tensor $S\in\binom{p}{q}$, and 
			another tensor $T\in\binom{r}{s}$. 

			Individually, the tensors act like so
			\[
				S[\theta_1,\cdots,\theta_p,\phi_1,\cdots,\phi_q]
			\]
			and
			\[
				T[\theta_1,\cdots,\theta_r,\phi_1,\cdots,\phi_s]
			\]
			The tensor product
			between them is
			\[
				S\otimes T \in \binom{p+r}{q+s}
			\]
			That is, it accepts $p+r$ covectors, and $q+s$ vectors. Let $\theta_i$ be covectors and 
			$\phi_j$ be vectors. 
			Then, the tensor product acts like
			\[
				(S\otimes T)[\theta_1,\cdots, \theta_{p},\phi_1,\cdots,\phi_q,\theta_{p+1},\cdots,\theta_{p+r},\phi_{q+1},\cdots,\phi_{q+s}]
			\]
			This can all be confusing, no? It is easy to make mistakes in this notation. As such, 
			you need to respect the order of the covectors and vectors. 
			We can rearrange the order by which the tensor accepts arguments. This is simply equivalent to
			\[
				(S\otimes T)[\theta_1,\cdots, \theta_{p+r}, \phi_1, \cdots, \phi_{q+s}]
			\]
			Because the order of the duals and vectors themselves do not matter, what matters is 
			the order of the vectors among the group of vectors, and the order of the covectors among the 
			group of covectors. That order is respected, but to simplify things, the new tensor acts the same way
			if we feed it the $S$ covectors first, then the $S$ vectors, followed by the $T$ covectors and vectors. 

			Writing it down in component form, we have
			\begin{align*}
				&S\otimes T =
				\tensor{S}{^{i_1,\cdots,i_p}_{j_1,\dots,j_q}}
				\tensor{T}{^{k_1,\cdots,k_q}_{l_1,\dots,l_s}}\\
						   &\p_{i_1}\otimes\cdots\otimes\p_{i_p}\otimes dx^{j_1}\otimes\cdots\otimes dx^{j_q}
				\p_{k_1}\otimes\cdots\otimes\p_{k_q}\otimes dx^{l_1}\otimes\cdots\otimes dx^{l_s}
			\end{align*}
			Equivalently,
			\begin{align*}
				&S\otimes T =
				\tensor{S}{^{i_1,\cdots,i_p}_{j_1,\dots,j_q}}
				\tensor{T}{^{k_1,\cdots,k_q}_{l_1,\dots,l_s}}\\
						   &\p_{i_1}\otimes\cdots\p_{i_p}\otimes
				\p_{k_1}\otimes\cdots\p_{k_q}\otimes
				dx^{j_1}\otimes\cdots\otimes dx^{j_q}\otimes
				dx^{l_1}\otimes\cdots\otimes dx^{l_s}
			\end{align*}
			We do this because the tensor product commutes for vectors and covectors.
			The slots for covectors and the slots for vectors should not be confused with 
			one another. Covector order and vector order matters, but the order of the covectors and vectors themselves don't really matter.

			Under this notation, there is no ambiguity. As long as the order of the covectors and vectors 
			are respected, we can swap them and indicate which slots belong to which.

			However, as you might have noticed, this notation is very cumbersome. I am intentionally
			showing you how inconvenient it is to use this notation. All those shit regarding 
			bases that you need to follow. However, mathematicians love this notation.
			They prefer it precisely because it is rigourous, and no confusion can happen.
			
			This notation is called the \textit{index-free} notation. It is index free because the 
			tensor itself has no index, only the component has. It is what we have been
			using so far to understand tensors. However, it has a few drawbacks.

			First, is that the expression of a tensor is very long and very cumbersome to write. It is
			easy to make mistakes. Second, is that you do not know the type of the tensor immediately just
			by looking at it. We need to specify that a tensor is of type $(k,l)$, which is separate
			from how we define the tensor itself. It is difficult to perform calculus on it, and 
			it's just inconvenient. Contraction operations on tensors using this notation is hard.
			Manipulating things is diffilcult. For example, showing that
			\[
				S\otimes T \neq T\otimes S
			\]
			explicity is cumbersome, since we need to define everything.

			However, this is very common in pure mathematics. It is inconvenient, but rigourous. It is unambigous,
			and it is exact.
		\item Index notation

			This notation is quite old fashioned. It is used by 
			physicists to express and  manipulate tensors, specially from those
			under the old universities such as Princeton or MIT. It was widely
			accepted in physics in the last century. Here, we denote the tensor with indices.

			For example, a $(2,1)$ tensor is just 
			\[
				\tensor{T}{^{\mu\nu}_{\sigma}}
			\]

			That's the tensor itself. All the basis that accompany it, and its components, are implied. 

			A tensor product between, say $\tensor{S}{^\mu_\nu}$ and $\tensor{T}{^{\alpha\beta^\gamma}_{\epsilon\xi}}$
			is just
			\[
				\tensor{S}{^\mu_\nu}
				\otimes
				\tensor{T}{^{\alpha\beta\gamma}_{\epsilon\xi}}
				=
				\tensor{(S\otimes T)}{^{\mu\alpha\beta\gamma}_{\nu\epsilon\xi}}
			\]

			It has a few advantages. First, you already know the type of the tensor without being explicitly
			told. From the symbol itself, for example, we can already tell that the tensor above is a $(4,3)$ tensor.
			Second, manipulating it for calculus, contraction, and other operations are easy. There is no need
			to fuss about with the basis vectors or covectors. It is convenient. Physicists are lazy, 
			we love our shortcuts, and we love manipulating numbers easily, and as such, while mathematicians 
			abhor this, physicists prefer it

			However, it also has some drawbacks. You can't easily tell if what you are manipulating 
			is the tensor itself, or just its components. It's hard to check the validity of the mathematics
			you do if it is difficult to identify what is it you are working with. Also, if you 
			manipulate larger tensors, the amount of indices become too many, and it might be 
			more difficult to work with. Lastly,
			mathematicians will not understand what you are referring to. If you give a mathametician, say
			\[
				\tensor{T}{^\mu_\beta}
			\]
			They will not call that a tensor. They will call that a scalar. It's the operation of a tensor
			on the bases. It is a collection of components.
			\[
				T[dx^\beta,\p_\mu]
			\]
			So for them, what physicists refer to as a tensor is not necesarilly a tensor itself. 
			So there is some difficulty collaborating between the two.

		\item Abstract Index notation

			Penrose and others, specially from those that came from the Chicago school of thought,
			went halfway between the two. They used both elements of the notation.

			They \textit{defined} the following: If a tensor symbol uses \textit{Latin} indices,
			it is a tensor, the geometric object.
			
			For example,
			\[
				\tensor{T}{^{a}_{bc}}
			\]
			From this, it is immediately clear that the tensor is a $(1,2)$ tensor. There is no
			explicit reference to its type, because it is obvious from the indices alone.

			A tensor product is simply
			\[
				\tensor{T}{^{ab}_{c}}
				\tensor{S}{^d_e} 
				=
				These constants describe the conserved quantities for a metric.
				\tensor{S}{^d_e}\tensor{T}{^{ab}_c}
				=
				\tensor{(S\otimes T)}{^{abc}_{de}} 
				\neq 
				\tensor{S}{^a_c}\tensor{T}{^{bd}_e}
			\]
			The fact that the tensor product is noncommutative is embedded in the notation.
			Since a tensor product makes a new tensor, it's also very easy to show.
			\[
				\tensor{T}{^{ab}_{c}}
				\tensor{S}{^d_e} 
				=
				\tensor{S}{^d_e}\tensor{T}{^{ab}_c}
				=
				\tensor{U}{^{abc}_{de}}
			\]
			From this, it is immediately clear that $U$ is a $(3,2)$ tensor. 

			It does not show the components, but it presents tensors in a more convenient way, similar 
			to the index notation that old physicists used.

			They then \textit{defined} that if a tensor symbol uses \textit{Greek} indices, then
			it is the components of a tensor, not the tensor itself. For example
			\[
				\tensor{T}{^\mu_{\nu\sigma}}
			\]
			are the tensor components. You can then combine it with the basis $p_\mu$, $dx^\nu$, and 
			others, so that it qualifies as a bona fide tensor.
			\[
				\tensor{T}{^\mu_{\nu\sigma}}\p_\mu \otimes dx^\nu \otimes dx^\sigma
			\]
			Therefore, a tensor can be represented as 
			\[
				\tensor{T}{^{a}_{bc}}
				=
				\tensor{T}{^\mu_{\nu\sigma}}\p_\mu\otimes dx^\nu\otimes dx^\sigma
			\]
			And a tensor acting on covectors and vectors can be easily represented as 
			\[
				T[\omega,\xi]=
				\tensor{T}{^{ab}_c}\omega_a \xi_b
			\]
			It is then immediately clear that $\omega$ and $\xi$ are covectors, and $T$ is
			lacking a vector input, since it is a $(2,1)$ tensor. 

			In addition, we can also immediately identify vectors as 
			\[
				v^a
			\]
			and covectors
			\[
				\omega_b
			\]
			You do not immediately
			get that information in index-free notation, most have to be explicitly stated, and you
			have to rely on context to understand things (that's why we had to do the overline underline notation to keep track of things earlier). It is also clear that it is a tensor \textit{acting} on other tensors,
			and not just components, and the indices indicate the slots being filled. 
			There is no ambiguity, but it's easy to use. 

			You therefore get both the convenience
			of index notation, and the rigour (should you desire) of the index-free notation.
			But tensor manipulation and operations itself are easier than the latter, but clearer than the former.

			I\footnote{Sir Vega referring to himself} myself trained 
			as a relativist and took my PhD and Post Doc under those schools, my adviser studied
			under Penrose, so I also adapt this notation.

			One major drawback of this is again, if you work with higher dimensional tensors,
			you will run out of letters very quickly. There have been instances where physicists 
			have resorted to using Hebrew and Japanese symbols to keep track of their tensors.
	\end{enumerate}

	Moving forward, we will use this abstract index notation to do our calculations, while sometimes still using index-free.
	You may ask, ``Why teach index free notation at all''? It's that you can understand
	all the rigour of the mathematics before you can be free to break it and make heuristics 
	of your own. Index free notation is cumbersome, but it is exact, which is useful.

	Furthermore, if I did not teach you index-free notation, you would have a hard time understanding
	mathematical papers and books teaching differential geometry from a pure math perspective. 
	It would take some time and difficulty before you fully grasp what they mean, and there are very good, 
	comprehensive pure math differential geometry text books that are written purely under index-free
	notation. 

	You would not gain anything from those resources if you do not understand the notation they use.
	We learned it first, so that we know the foundations of what really it is we are doing.
	When we have truly leard all the rules, then we can start being lazy and do things
	in the abstract way we desire. However, never forget that the root of the notation you are using,
	which is the index-free rigourous definition. If you take your notation and run with it, 
	without appreciating the true meaning underneath it all, you will have a hard time extending
	your mathematics beyond where the assumptions your notation makes hold.
